% ------------------------------------------------------------------------
% bjourdoc.tex for birkjour.cls*******************************************
% ------------------------------------------------------------------------
%%%%%%%%%%%%%%%%%%%%%%%%%%%%%%%%%%%%%%%%%%%%%%%%%%%%%%%%%%%%%%%%%%%%%%%%%%

\documentclass{birkjour}
%
%
% THEOREM Environments (Examples)-----------------------------------------
%
 \newtheorem{thm}{Theorem}[section]
% \newtheorem{cor}[thm]{Corollary}
% \newtheorem{lem}[thm]{Lemma}
% \newtheorem{prop}[thm]{Proposition}
% \theoremstyle{definition}
 \newtheorem{defn}[thm]{Definition}
% \theoremstyle{remark}
% \newtheorem{rem}[thm]{Remark}
% \newtheorem*{ex}{Example}
 \numberwithin{equation}{section}

\usepackage[noadjust]{cite}
\usepackage{amsfonts}
\usepackage{listings}
\usepackage{algorithm}
\usepackage{algorithmic}
\usepackage{booktabs}
\usepackage{float}
\usepackage{caption}
\usepackage{tikz}
\usepackage{hyperref}
\begin{document}

%-------------------------------------------------------------------------
% editorial commands: to be inserted by the editorial office
%
%\firstpage{1} \volume{228} \Copyrightyear{2004} \DOI{003-0001}
%
%
%\seriesextra{Just an add-on}
%\seriesextraline{This is the Concrete Title of this Book\br H.E. R and S.T.C. W, Eds.}
%
% for journals:
%
%\firstpage{1}
%\issuenumber{1}
%\Volumeandyear{1 (2004)}
%\Copyrightyear{2004}
%\DOI{003-xxxx-y}
%\Signet
%\commby{inhouse}
%\submitted{March 14, 2003}
%\received{March 16, 2000}
%\revised{June 1, 2000}
%\accepted{July 22, 2000}
%
%
%
%---------------------------------------------------------------------------
%Insert here the title, affiliations and abstract:
%


\title[Parallel Transport of Vectors on 3D Mesh by the Vector Heat Method using FEM]
 {Parallel Transport of Vectors on 3D Mesh by the Vector Heat Method using FEM}

%----------Author 1
\author[Mauricio Cele Lopez Belon]{Mauricio Cele Lopez Belon}
\address{Madrid, Spain}
\email{mclopez@outlook.com}

%----------classification, keywords, date
\subjclass{Parallel algorithms 68W10; Clifford algebras, spinors 15A66}

\keywords{Geometric Algebra, Quaternion Estimation, Mesh Deformation}

\date{October 21, 2024}
%----------additions
%\dedicatory{To my wife}
%%% ----------------------------------------------------------------------

\begin{abstract}

Compute the parallel transport of tangent vectors defined on a 3D mesh by the Vector Heat Method using Finite Element Method (FEM).

\end{abstract}

%%% ----------------------------------------------------------------------
\maketitle
%%% ----------------------------------------------------------------------
%\tableofcontents
\section{Introduction}

This follows the method described in \cite{Sayas2015} and \cite{Langtangen2014}.

Solve the following Heat Equation PDE with Dirichlet boundary conditions (boundary value problem):

\begin{tabular}{l l}
	Solve $\nabla^2 U(\vec X, t) = \frac{ \partial U(\vec X, t)}{\partial t}$ & in region $\Omega$ \\
	s.t. $U(\vec X, t) = G(\vec X)$  &  on the region boundary $\partial \Omega$ for all times \\
	$U(\vec X, 0) = U_0(\vec X)$  &  initial condition at time t = 0 \\
\end{tabular}
where $U(\vec X, t)$ is a vector field and $\nabla^2$ is the ``Connection Laplacian'' which produces another
vector field ans is defined in terms of the covariant derivative $\nabla$ as:

\begin{eqnarray}
	\label{eqn:connection_laplacian}
	\nabla^2 U(\vec X) = \nabla^* \nabla U(\vec X) 
\end{eqnarray}
where $\nabla^*$ is the adjoint of the covariant derivative $\nabla$, $U(\vec X)$ is a tangent vector field on the 3D surface 
(2D manifold embedded in $\mathbb R^3$). $G(\vec X)$ the vector boundary conditions and $U_0(\vec X)$ is a vector field at time 0.

The covariante derivative $\nabla(U, w)$ is a ``directional derivative'' operator that has two vector field arguments:
$U(\vec X)$ and $w(\vec X)$. Where $U(\vec X)$ is the vector field which is differentiated and $w(\vec X)$ is a vector field 
telling the direction in which directional derivatives are taken. It returns a vector field as which is a measure of 
relative change of $U$ w.r.t to $w$. 

Usually one can omit the vector field $w$ in which case we talk about a total derivative $\nabla U$ which is a 2-tensor 
(i.e., a matrix). For example in flat space $\mathbb R^3$ the vector field represented as $U = \left[u_x, u_y, u_z\right]^T$ 
the operator $\nabla U$ would be:

\begin{eqnarray}
	\nabla U =  
	\left[\begin{array}{c}
		\frac{\partial}{\partial x} \\
		\frac{\partial}{\partial y} \\
		\frac{\partial}{\partial z} 
	\end{array}\right] 
	\left[\begin{array}{ccc}
		u_x &  u_y  & u_z 
	\end{array}\right] 
    =
	\left[\begin{array}{ccc}
		\frac{\partial}{\partial x} u_x  &  \frac{\partial}{\partial y} u_x   &   \frac{\partial}{\partial z} u_x \\
		\frac{\partial}{\partial x} u_y  &  \frac{\partial}{\partial y} u_y   &   \frac{\partial}{\partial z} u_y \\
		\frac{\partial}{\partial x} u_z  &  \frac{\partial}{\partial y} u_z   &   \frac{\partial}{\partial z} u_z 
	\end{array}\right] \nonumber
\end{eqnarray}

In more compact notation
\begin{eqnarray}
	\nabla U = \partial \otimes U = J(U)
\end{eqnarray}

where $\otimes$ is the ``tensor product'', $J(U)$ is the Jacobian matrix of $U$ (or the differential) and 
$\partial$ is the ``grandient operator''. 

In a curved manifold the above definition would need to be adjusted to cope with curvature with the aid of a ``connection''.

The adjoint of $\nabla$ can be defined by ``integration by parts'' 
(see  \href{https://en.wikipedia.org/wiki/Vector_calculus_identities}{wikipedia}) 
over surface $\Omega$:

Recall adjoint operator is defined as an inner product:
\begin{eqnarray}
	\left\langle f , \nabla^* g \right\rangle = \left\langle \nabla f , g \right\rangle \nonumber\\
	\int { f \ \nabla^* g \ dx }  = \int { \nabla f \ g \ dx } \nonumber
\end{eqnarray}
recall intergation by parts formula:
\begin{eqnarray}
	\int_{\Omega} { f'(x) g(x) \ dx }  = \int_{\partial \Omega} { f(x) g(x) \ dx } - \int_{\Omega} { f(x) g'(x) \ dx } \nonumber
\end{eqnarray}
doing integration by parts on $\int { \nabla_x f \ g \ dx }$ we get:
\begin{eqnarray}
	\int_{\Omega} { \nabla_X f \ g \ dx }  = \int_{\partial \Omega} { f \ g \ dx } - \int_{\Omega} { f \ \nabla_X^* g \ dx } \nonumber
\end{eqnarray}
where $\nabla_X$ is the directional derivative such that $\nabla_X f = X \cdot \nabla f$

Applying the above to the vector field valued functions $F$ and $G$
\begin{eqnarray}
	\label{eqn:integration_by_parts_adjoint}
	\int_{\Omega}{ \nabla_X F \cdot G \ dx } = \int_{\Omega}{ F \cdot \nabla_X^* G \ dx } + \int_{\partial \Omega}{ F \cdot G(N) \ dS}
\end{eqnarray}
where $\nabla_X F$ and $\nabla_X^* G$ are vector fields and the $(\cdot)$ is inner product.
$N$ is unit normal at boundary $\partial \Omega$. For details check 
\href{https://ddg.math.uni-goettingen.de/pub/vector-dirichlet.pdf}{A Simple Discretization of the Vector Dirichlet Energy}
For identities in flat space check \href{https://thatsmaths.com/2020/01/23/adjoints-of-vector-operators/}{adjoints of vector operators}

Going back to the original goal it is to compute the parallel transported vector field $U$ from vector fields
given at boundary. The Vector Heat Method works as follows: solves the Vector Heat Equation 
PDE $\frac{\partial U}{\partial t} = \nabla^2 U$ with dirichlet boundary conditions using Finite Element Method (FEM)

\begin{tabular}{l l}
	Solve $\bigtriangleup U(\vec X, t) = \frac{ \partial U(\vec X, t)}{\partial t}$ & in region $\Omega$ \\
	s.t. $U(\vec X, t) = G(\vec X)$  &  on the region boundary $\partial \Omega$ for all times \\
	$U(\vec X, 0) = U_0(\vec X)$  &  initial condition at time t = 0 \\
\end{tabular}

\subsection{The Vector Heat PDE}

Since $\frac{ \partial U(\vec X, t)}{\partial t}$ is a time derivative we discretize the time using Backguard Euler
\begin{eqnarray}
	\frac{ \partial U(\vec X, t)}{\partial t} = \frac{U(\vec X, t + \Delta t) - U(\vec X, t)}{\Delta t}
\end{eqnarray}

Replacing on original PDE for backguard euler:
\begin{eqnarray}
	\frac{U(\vec X, t + \Delta t) - U(\vec X, t)}{\Delta t} = \nabla^2 U(\vec X, t + \Delta t)
\end{eqnarray}

Re-arranging:
\begin{eqnarray}
 	U(\vec X, t + \Delta t) - \Delta t  \nabla^2 U(\vec X, t + \Delta t) = U(\vec X, t)
\end{eqnarray}

or better:
\begin{eqnarray} 
	U(\vec X, t_i) - \Delta t \nabla^2 U(\vec X, t_i) = U(\vec X, t_{i-1})
\end{eqnarray}

The integral form (notice $\int_{\Omega}$ is a double integral) is integrating each vector component separately.
\begin{eqnarray} 
	\int_{\Omega}{ U(\vec X, t_i) \ dA} - \Delta t \int_{\Omega}{ \nabla^2 U(\vec X, t_i) \ dA } = \int_{\Omega} { U(\vec X, t_{i-1}) \ dA }
\end{eqnarray}

\section{Weak Form}


The weak form comes from the ``variational method''. Here one tries to get a solution to the 
PDE by translating it into a minimization problem. That's done by definig an energy function that is a 
linear functional on a function space $\Omega$. The linear functional is artificial i.e., pure 
mathematical device to enable optimization. In consequence, solving the variational problem may 
require less continuity than that of the actual solution (hence the name WEAK form).
 
Consider the test function $\vartheta (x,y)$ (a.k.a. test function) with compact support defined on $\Omega$.
The ``Galerkin'' variational method states that the inner product of $\vartheta(x,y)$ with a residual function $R(x,y)$ should be zero.
Note that in this case $\vartheta$ have to be a vector field of the form 
$\vartheta(\vec X) = \left[\varphi_i(\vec X), \varphi_j(\vec X)\right]$ which is 
a vector with a scalar test functions with compact support as components.

\begin{eqnarray}
	\int_{\Omega} { \vartheta(\vec X) \cdot R(\vec X) \ dA } = 0 \nonumber
\end{eqnarray}
 where the $(\cdot)$ is inner product and $R(x,y)$ is the residual or error of the function i.e., 
\begin{eqnarray}
	R(x,y) = U(\vec X, t_i) - \Delta t \nabla^2 U(\vec X, t_i) - U(\vec X, t_{i-1}) \nonumber
\end{eqnarray}

The weak form of the PDE is:
\begin{eqnarray}
	\vartheta(\vec X) \cdot R(\vec X) = 0 \nonumber\\
	\vartheta(\vec X) \cdot (U(\vec X, t_i) - \Delta t \nabla^2 U(\vec X, t_i) - U(\vec X, t_{i-1})) = 0 \nonumber\\
	U(\vec X, t_i) \cdot \vartheta(\vec X) - \Delta t \vartheta(\vec X) \cdot \nabla^2 U(\vec X, t_i) = U(\vec X, t_{i-1}) \cdot \vartheta(\vec X) \nonumber
\end{eqnarray}
 and the integral (variational) form of the PDE is
\begin{eqnarray}
	\label{eqn:original_weak_form}
	\int_{\Omega}{ U_t \cdot \vartheta \ dA } - \Delta t \int_{\Omega}{ \vartheta \cdot \nabla^2 U_t \ dA} = \int_{\Omega} { U_{t-1} \cdot \vartheta \ dA }
\end{eqnarray}
NOTE on inner products: The integral form is more important since integral of
of a ``vector field'' function $f(x,y)$ times test function $\vartheta(x,y)$ is the ``inner product'' 
$\int_{\Omega}{ f(x,y) \cdot \vartheta(x,y) \ dA }$ where the $(\cdot)$ is inner product of vectors.
If function $\vartheta(x,y)$ is a ``test function" which is a vector field with separate scalar test function with compact 
support as components then the inner product is way to ``measure'' the 
function $f(x,y)$ on location where $\vartheta(x,y)$ is defined.

Replacing $F = \vartheta$ and $G = \nabla U$ in (\ref{eqn:integration_by_parts_adjoint}) we get:

\begin{eqnarray}
	\label{eqn:covariant_derivative_adjoint}
	\int_{\Omega}{ \nabla \vartheta : \nabla U \ dx } = \int_{\Omega}{ \vartheta \cdot \nabla^* \nabla U \ dx } + \int_{\partial \Omega}{ \vartheta \cdot \nabla U(N) \ dS} \nonumber\\
	\int_{\Omega}{ \vartheta \cdot \nabla^* \nabla U \ dx } = \int_{\Omega}{ \nabla \vartheta : \nabla U \ dx } - \int_{\partial \Omega}{ \vartheta \cdot \nabla_N U \ dS} \nonumber\\
	\int_{\Omega}{ \vartheta \cdot \nabla^2 U \ dx } = \int_{\Omega}{ \nabla \vartheta : \nabla U \ dx } - \int_{\partial \Omega}{ \vartheta \cdot \nabla_N U \ dS}
\end{eqnarray}
where $\nabla U$ and $\nabla \vartheta$ are rank-2 tensor fields by virtue of the covariant derivative. 
Setting $\vartheta = 0$ at the boundary and replacing (\ref{eqn:covariant_derivative_adjoint}) in (\ref{eqn:original_weak_form})

\begin{eqnarray}
	\label{eqn:weak_form}
	\int_{\Omega}{ U_t \cdot \vartheta \ dA } - \Delta t \int_{\Omega}{ \nabla \vartheta : \nabla U_t \ dA } = \int_{\Omega} { U_{t-1} \cdot \vartheta \ dA }
\end{eqnarray}


\section{Discretization in Flat Space $\mathbb R^2$}
\label{sec:discretization} 

The domain $\Omega$ has to be meshed in $N$ nodes and $T$ triangles. For linear elements the nodes coincides with 
mesh vertices (but for quadratic elements there will be nodes defined on the triangle's edge mid-points, 
so there can be more nodes than vertices).

The ``test function'' is a function that allows us to linearly interpolate local orthonormal frames 
located at mesh nodes to the triangle-set sorrounding the node. Those local orthonormal frames 
located at mesh nodes are ``tangent spaces'' of the nodes. In the flat space of $\mathbb R^2$ all
tangent spaces are equal, so the interpolation of them is simple. However the key point to understand 
is that test functions represent linear interpolation (and atenuation) of ``local frames''. 
We define a test function $\varphi_i$ per node $N_i$ which will be used to linearly interpolate tangent vectors
defined in terms of local frames localted at nodes to all triangles incident to the node $N_i$. 


Since we are in flat space $\mathbb R^2$, the tangent space at node $N_i$ is the same as any other tangent space everywhere else 
in the domain $\Omega$. Therefore parallel transport is just the identify function or ``teleportation'', 
and the covariant derivative $\nabla_X U$ coincides with the directional derivative $X \cdot \nabla U$.

The ``local basis'' for tanget space at $N_v$ is defined as $T_v = (e_1, e_2)$, which is an orthonormal frame. 
The ``local basis'' for tanget space at triangle is $T_f = (e_1, e_2)$, the connection 1-form $\omega_i^f = I$ where $I$ is 
the identity matrix. Therefore the vectors at $U(x,y)$ are expressed in terms of the ``global basis``'' $(e_1, e_2)$

As usual, we define a set of scalar tests functions with compact support $\phi_i(x, y)$ 
for every node $N_i$ to be equal to $0$ or $1$ and linear everywhere else:
\begin{eqnarray}
	\phi_j(\vec N_j) = 1 \ \ \  \forall j = i \nonumber\\
	\phi_j(\vec N_j) = 0 \ \ \  \forall j \neq i \nonumber\\
	\phi_j(\vec X) = \text{linear in} \ \Omega_i \nonumber
\end{eqnarray}

We use those $\phi_j$ to define the vector-valued basis function $\varphi_i$ (see 
\href{https://math.stackexchange.com/questions/1320672/system-of-equations-for-vector-valued-functions-problems}{ref1}, 
\href{https://www.dealii.org/developer/doxygen/deal.II/group__vector__valued.html}{ref2}, 
\href{https://www.dealii.org/developer/doxygen/deal.II/step_8.htmls}{ref3},
\href{https://scicomp.stackexchange.com/questions/40735/fem-for-vector-valued-problems-reference-request}{ref4}):

\begin{eqnarray}
	\varphi_i(\vec X) = \vec \phi_i(\vec X) T_j \nonumber\\
	\varphi_i(\vec X) = 
		\left[\begin{array}{c}
			\phi_i(\vec X) \\
			\phi_i(\vec X)
		\end{array}\right]^T (e_1, e_2)
\end{eqnarray}
Where $T_j = (e_1, e_2)$ is the ``local basis'' ($2 \times 2$ matrix) for tanget space at $N_j$, 
$\vec X$ is vector expressed also in basis $(e_1, e_2)$.

Now we discretise the solution $U(x,y)$ to be a finite sum:
\begin{eqnarray}
	U(x,y) = \sum_j^N{ \varphi_j(x,y)^T
	\left[\begin{array}{cc}
		U^{e_1}_j & 0\\
		0 & U^{e_2}_j
	\end{array}\right]}
\end{eqnarray} 
where $\odot$ is the component-wise product of vectors 
(a.k.a., \href{https://en.wikipedia.org/wiki/Hadamard_product_%28matrices%29}{Hadamard product})
and $U_j$ is a vector ($2 \times 1$ column).

Now we discretise $\nabla U(x,y)$ which is a rank-2 tensor field
\begin{eqnarray}
	\nabla U(x,y) = \sum_j^N{ \nabla \varphi_j(x,y) U_j } \nonumber\\
	\nabla U(x,y) = \sum_j^N{ 
	\left[\begin{array}{c}
		\nabla \phi_j(\vec X) \\
		\nabla \phi_j(\vec X)
	\end{array}\right]^T 
	\left[\begin{array}{cc}
		U^{e_1}_j & 0\\
		0 & U^{e_2}_j
	\end{array}\right] }
\end{eqnarray} 
where $\nabla \varphi_j$ is a $2 \times 2$ matrix.

Given $\vartheta(x,y) = \left[\phi_i(x,y), \phi_i(x,y)\right]^T$ ,
we discretise the test function $\nabla \vartheta(x,y)$ in a the following manner:
\begin{eqnarray}
	\nabla \vartheta(x,y) = \sum_j^N{ \nabla \psi_j(x,y) }  \nonumber\\
	\nabla \vartheta(x,y) = \sum_j^N{ 
	\left[\begin{array}{c}
		\nabla \phi_j(\vec X) \\
		\nabla \phi_j(\vec X)
	\end{array}\right] }
\end{eqnarray} 

We discretise the $\nabla \vartheta : \nabla U$:
\begin{eqnarray}
	\nabla \vartheta : \nabla U = \left( \sum_i{ \nabla \psi_i }\right) : \left( \sum_j{ \nabla \varphi_j U_j } \right) \\
	\nabla \vartheta : \nabla U = \sum_{ij}{ \left( \nabla \psi_i : \nabla \varphi_j \right) U_j } \nonumber\\
	\nabla \vartheta : \nabla U = \sum_{ij}{ Tr\left( 
		\left[\begin{array}{cc}
			\nabla_{e_1} \phi_i(\vec X) & \nabla_{e_2} \phi_i(\vec X)\\
			\nabla_{e_1} \phi_i(\vec X) & \nabla_{e_2} \phi_i(\vec X)
		\end{array}\right]^T
		\left[\begin{array}{cc}
			\nabla_{e_1} \phi_j(\vec X) & \nabla_{e_2} \phi_j(\vec X)\\
			\nabla_{e_1} \phi_j(\vec X) & \nabla_{e_2} \phi_j(\vec X)
		\end{array}\right] U_j
		\right) } \nonumber\\
	\nabla \vartheta : \nabla U = \sum_k \sum_{ij}{ \left( \nabla \phi_i(\vec X) \cdot \nabla \phi_j(\vec X)\right) U^{e_k}_j } \nonumber
\end{eqnarray}
where the contraction of the two rank-2 tensors is another rank-2 tensor e.g., $G_i \cdot G_j = G_i^T G_j$ or in tensor notation
$G_{ij} \ G_{jk} = G_{ik}$. The integrals are now much smaller i.e., per triangle-set $\Omega_i$ around nodes 
$\vec N_i$ or tent function $\varphi_i(\vec X)$

The stiffness matrix for triangle $T$ is:
\begin{eqnarray}
	M_k(i,j) = I_{2 \times 2} \int_{\Omega}{ \nabla \phi_i \cdot \nabla \phi_j \ dA }
\end{eqnarray}

So we can form the linear system:
\begin{eqnarray}
	\sum_{ij}{ I_{2 \times 2} \left( \int_{\Omega}{ \nabla \phi_i \cdot \nabla \phi_j \ dA }\right) U_j }
\end{eqnarray}
where $U_j$ is $2 \times 1$ column vector and $I_{2 \times 2}$ is identity matrix.


\section{Assembly}


The process of creating the stiffness matrix $M$ is called assembly. It consist of creating the matrix $M_k$ for each 
triangle $T_k$ and then add it to the global $M$ matrix. 
 
The global matrix $M$ is $2 N \times 2 N$  where $N$ is number of mesh nodes. In previous section~\ref{sec:discretization}
the matrix $M_k$ is a $2 N \times 2 N$ a super-sparse block matrix where each entry is $2 \times 2$ block matrix just with
a few non-zero entries which corresponds to a single triangle contribution.
 \begin{eqnarray} 
	M = \sum_k M_k
\end{eqnarray}
 
In practice the $M_k$ matrix is not $2 N \times 2 N$ super-sparse, but a $6 \times 6$ matrix (in case of linear triangle elements).
There is a map from indices of the $6 \times 6$ matrix to the $2 N \times 2 N$ matrix as follows.
 
Each triangle's nodes has a global index assigned in the mesh, we also establish a local indexing of triangle
like $1$, $2$ and $3$. We define a trivial bijective map between global and local indices.


\section{Discretization in 3D curved space e.g., 3D surface}

The 3D curved domain $\Omega$ has to be discretised in $N$ nodes and $T$ triangles. 
For linear elements the nodes coincides with vertices, but for 
quadratic elements there will nodes defined on the mid-point of edges of triangles, so there will be more nodes than vertices.

We define a test function $\Phi_i$ per node $N_i$ which will be used to linearly interpolate basis tangent vectors located at
mesh nodes to all triangles incident to the node $N_i$. 

First we define a local basis for tangent space $T_v$ per node $N_i$ in which tangent vector of $U(x,y)$ will reside. We will also 
need to define a local tangent space basis per each triangle $T_f$, which not coincide with $T_v$, that will be used to 
interpolate basis tanegt vectors within triangles. Finally we need a ``connection'' $R_i^f$ which parallel transports a basis 
of tangent vectors $T_v$ at node $N_i$ to a basis of tangent vectors $T_f$ at triangle $T_f$ incident to node $N_i$:

\begin{eqnarray}
	\begin{array}{cccc}
		R_i^f =   & T_f^T     &  Q_{vf}    &   T_v       \\
		2\times2  & 2\times3  &  3\times3  &   3\times2		 
	\end{array}
\end{eqnarray}
where:

\begin{tabular}{p{0.05\linewidth} p{0.80\linewidth}}
	$R_i^f$ & is $2\times2$ matrix connecting 2 tangent spaces: the local tangent space at node $N_i$ to a local tangent space 
	at triangle $T_f$. The tangent spaces are connected in such a way that a tangent vector is ``parallel transported'' from
	one tangent space to the next \\
	$T_f$ & is the $3\times2$ local tangent space basis for triangle $f$: $T_f = \left[e_i, e_j\right]^T$ \\
	$Q_{vf}$ & is the $3\times3$ rotation matrix aligning node $N_i$ tangent plane with face $T_f$ plane. \\
	$T_v$ & is the $3\times2$ local tangent space basis for tangent space at node $N_i$: $T_v = \left[e_i, e_j\right]^T$
\end{tabular}

The connection $R_i^f$ defines the mapping $\Phi_i^f$ at node $N_i$ restricted to triangle $f$:
\begin{eqnarray}
	\Phi_i^f(\vec X) = R_i^f \vec X
\end{eqnarray}
which parallel transport a tangent vector $\vec X$ in tangent space basis at node $N_i$ to a triangle $T_f$ local basis.

We also define scalar ``tests'' functions with compact support $\phi_i(x, y)$ for every node $N_i$ to be equal to $0$ or $1$
and linear everywhere else:
\begin{eqnarray}
	\phi_j(\vec N_j) = 1 \ \ \  \forall j = i \nonumber\\
	\phi_j(\vec N_j) = 0 \ \ \  \forall j \neq i \nonumber\\
	\phi_j(\vec X) = \text{linear in} \ \Omega_i \nonumber
\end{eqnarray}

We use those $\phi_j$ to define the vector-valued basis function $\Psi_i^f$:
\begin{eqnarray}
	\Psi_i^f(x,y)(\vec X) = \phi_j(x,y) \Phi_i^f \vec X \nonumber\\
	\Psi_i^f(x,y)(\vec X) = \phi_j(x,y) R_i^f \vec X
\end{eqnarray}

Now we discretise the solution $U(x,y)$ on triangle $T_f$ to be a finite sum:
\begin{eqnarray}
	U(x,y) = \sum_j{ \Psi_j(x,y) U_j } \nonumber\\
	U(x,y) = \sum_j{ \phi_j(x,y) R_j^f U_j }
\end{eqnarray}
and the covariant derivative $\nabla U(x,y)$:
\begin{eqnarray}
	\nabla U(x,y) = \sum_j{ \nabla \Psi_j(x,y) U_j } \nonumber\\
	\nabla U(x,y) = \sum_j{ \nabla \phi_j(x,y) \otimes R_j^f U_j } \nonumber\\
\end{eqnarray}
or in matrix form:
\begin{eqnarray}
	\nabla U(x,y) = \sum_j{ R_j^f 
	\left[\begin{array}{cc}
		U^{e_1}_j & 0\\
		0 & U^{e_2}_j
	\end{array}\right]
	\left[\begin{array}{c}
		\nabla \phi_j(x,y) \\
		\nabla \phi_j(x,y)
	\end{array}\right]}
\end{eqnarray}
which is a $2\times2$ matrix.

We discretise the test function $\vartheta(x, y)$ on triangle $T_f$ as:
\begin{eqnarray}
	\vartheta(x, y) = \sum_j{ R_j^f \psi_j(x,y) } \nonumber\\
	\vartheta(x, y) = \sum_j{ R_j^f \left[\phi_j(x,y), \phi_j(x,y)\right]^T }
\end{eqnarray}
where $\psi_j(x,y) = \left[\phi_j(x,y), \phi_j(x,y)\right]^T$ and its covariant derivative $\nabla \vartheta(x, y)$ as:
\begin{eqnarray}
	\nabla \vartheta(x, y) = \sum_j{ R_j^f \nabla \psi_j(x,y) } \nonumber\\
	\nabla \vartheta(x, y) = \sum_j{ R_j^f \left[\nabla \phi_j(x,y), \nabla \phi_j(x,y)\right]^T }
\end{eqnarray}

We discretise the $\nabla \vartheta : \nabla U$ for each element $T_f$ on the triangle-set around node $N_i$ 
as contribution of nodes $N_j$. Using the fact that $A : B = Tr(A^T B)$ and the cyclic property of the trace we get:
\begin{eqnarray}
	\nabla \vartheta : \nabla U = \left(\sum_i{ R_i \nabla \psi_i } \right) : \left( \sum_j{ \nabla \phi_j \otimes R_j U_j } \right) \nonumber\\
	\nabla \vartheta : \nabla U = \sum_{ij}{ R_i \nabla \psi_i : \nabla \phi_j \otimes R_j U_j } \nonumber\\
	\nabla \vartheta : \nabla U = \sum_{ij}{ Tr(\nabla \psi_i^T R_i^T R_j [U_j] \nabla \psi_j) } \nonumber\\
	\nabla \vartheta : \nabla U = \sum_{ij}{ Tr( \nabla \psi_j \nabla \psi_i^T R_i^T R_j [U_j] ) }
\end{eqnarray}
Where $[U_j] = 	\left[\begin{array}{cc}
	U^{e_1}_j & 0\\
	0 & U^{e_2}_j
\end{array}\right]$. Noting that 
$Tr(\nabla \psi_j \nabla \psi_i^T) = Tr(\nabla \psi_i^T \nabla \psi_j) = \nabla \phi_i \cdot \nabla \phi_j$:
\begin{eqnarray}
	\nabla \vartheta : \nabla U = \sum_k \sum_{ij}{ \left(\nabla \phi_i \cdot \nabla \phi_j\right) (R_i^T R_j U_j) \cdot e_k }
\end{eqnarray}

The stiffness matrix for triangle $T$ is a $6\times6$ matrix:
\begin{eqnarray}
	M_k(i,j) = R_i^T R_j \int_{\Omega}{ \nabla \phi_i \cdot \nabla \phi_j \ dA }
\end{eqnarray}

So we can form the linear system:
\begin{eqnarray}
	\sum_{ij}{ \left( \int_{\Omega}{ \nabla \phi_i \cdot \nabla \phi_j \ dA }\right) R_i^T R_j U_j }
\end{eqnarray}
where $U_j$ is $2 \times 1$ column vector and $R_i^T R_j$ is the discrete connection $2\times 2$ matrix.

\section{The Covariant Derivative $\nabla \Psi$}

We use the Cartan's moving frames formalism. Moving frames are orthonormal frames attached to each point on the surface. Those 
frames lies in the tangent space associated to each point, thus the frames are orthogonal unit tangent vectors. Those frames 
don't form a ``holonomic'' basis i.e., those are not vector fields associated to a directional derivative or partial derivative 
operator. In particular, moving frames are not continuous over the surface and its length is always $1$ (in contrast, holonomic 
basis can have different lengths, thus their inner product defines the metric tensor). For this reason the Crhistoffel Symbols 
is a concept associated to differentiation of holonomic basis while ``Connection Forms'' is a concept associated with 
``exterior differentiation'' of moving frames. However in the flat space $\mathbb R^n$ under cartesian coordinates the 
Christoffel Symbols and Connection Forms are the same.

First recap definition of covariant derivative of a vector field $U$ in direction of vector field $w$ 
(see \href{https://en.wikipedia.org/wiki/Connection_form#Example:_the_Levi-Civita_connection}{the Levi-Civita connection}):

\begin{tabular}{p{0.25\linewidth} p{0.65\linewidth}}
	$E = (e_1, e_2)$ & local frame field i.e., moving frame seen as pair of tangent vectors $e_1$ and $e_2$ \\
	$U = U_1 e_1 + U_2 e_2$ & vector field expressed in moving frames as $\sum_i U_i e_i$  \\
	$w = w_1 e_1 + w_2 e_2$ & vector field expressed in moving frames as $\sum_i w_i e_i$  \\
\end{tabular}
\begin{eqnarray}
	\omega(e_k) = \left[\begin{array}{cc}
		\omega_{11}(e_k) & \omega_{12}(e_k) \\
		\omega_{21}(e_k) & \omega_{22}(e_k)
	\end{array}\right] \text{Matrix of 1-forms a.k.a. connection form} \nonumber
\end{eqnarray}

We denote $\partial U_i$ the gradient of scalar component $U_i$ which is a vector 
$\partial U_i = \left[\frac{\partial}{\partial e_1} U_i, \frac{\partial}{\partial e_2} U_i\right]^T$. 
And $\partial U$ the Jacobian matrix of $U$
\begin{eqnarray}
	\partial U = \left[\partial U_1, \partial U_2\right]^T \nonumber\\
	\partial U =\left[\begin{array}{cc}
		\frac{\partial}{\partial e_1} U_1 & \frac{\partial}{\partial e_2} U_1 \\
		\frac{\partial}{\partial e_1} U_2 & \frac{\partial}{\partial e_2} U_2
	\end{array}\right]
\end{eqnarray}

The covariant derivative $\nabla_w U$ of a vector field $U$ in direction of vector field $w$ is defined as
\begin{eqnarray}
	\nabla_w U = w \cdot \nabla U \nonumber\\
	\nabla_w U = w_1 \nabla_{e_1} U + w_2 \nabla_{e_2} U 
\end{eqnarray}

Now we need to expand $\nabla_{e_i} U$:
\begin{eqnarray}
	\nabla_{e_j} U = \partial U e_j + \omega(e_j) U \nonumber\\
	\nabla_{e_j} U =
	\left[\begin{array}{c}
		\partial U_1 \cdot e_j + \omega_{11}(e_j) U_1 + \omega_{12}(e_j) U_2 \\
		\partial U_2 \cdot e_j + \omega_{21}(e_j) U_1 + \omega_{22}(e_j) U_2
	\end{array}\right] 
\end{eqnarray}
 Putting back $w = w_1 e_1 + w_2 e_2$ we get:

\begin{eqnarray}
	\nabla_w U = \partial U w  + \omega(w) U
\end{eqnarray}

The above expression coincides with definition in Liu, Tong, De Goes, Desbrun ``Discrete Connection and Covariant 
Derivative for Vector Field Analysis and Design'' where
$\partial U w$ is the multiplication of the Jacobian of $U$ with $w$ which is giving the directional derivative
in flat space.
$\omega(w) U$ is the multiplication of the connection form $\omega(w)$ with $U$ which is parallel transporting
$U$ to the tip of $w$.

We can define $\nabla U$ just in terms of $U$ without a particular direction $w$ as argument
(see \href{https://math.stackexchange.com/questions/3468914/covariant-derivative-of-a-vector-field}{covariant derivative of a vector field}).
\begin{eqnarray}
	\nabla_w U = \partial_{e_k} U_i w_k  + \omega_{ij}^k w_k U_i \nonumber\\
	\nabla_w U = (\partial_{e_k} U_i  + \omega_{ij}^k U_i) w_k \nonumber\\
	\nabla U = \partial_{e_k} U_i  + \omega_{ij}^k U_i
\end{eqnarray}
which can also be written as:
\begin{eqnarray}
	\nabla U = \partial U + \omega_{ij}^k U_i
\end{eqnarray}

With the above representation we can contract $\nabla U$ with vector $w$ to get the classic 
Levi-Civita connection:
\begin{eqnarray}
	\nabla_w U = w \cdot \nabla U
\end{eqnarray}

Recall the definition of $\Psi(x,y)$ basis function which is defined for every point $(x, y)$ on triangle $T_f$ 
applied to a tangent vector $\vec X_i$ which is defined at node $N_i$:
\begin{eqnarray}
	\Psi_i(x,y)(\vec X_i) = \phi_i(x,y) R_i^f \vec X_i
\end{eqnarray}
The basis function $\Psi$ does a parallel transport of vectors from tanget space at $N_i$ to a flat tangent space at triangle
$T_f$ by using a discrete connection $R_i^f$. The discrete connection $R_i^f$ is in essence the Connection Form $\omega(\vec X_i)$.
The covariant derivative $\nabla \Psi$ is simply:
\begin{eqnarray}
	\nabla \Psi_i(x,y)(\vec X_i) = \nabla \phi_i(x,y) \otimes R_i^f \vec X_i + \phi_i(x,y) \nabla (R_i^f \vec X_i) \nonumber\\
	\nabla \Psi_i(x,y)(\vec X_i) = \partial \phi_i(x,y) \otimes R_i^f \vec X_i
\end{eqnarray}
where $\partial \phi_i(x,y)$ is the grandient of a scalar valued function $\phi_i(x,y)$, which is a vector and $\otimes$ is the tensor
(or dyadic) product of vectors. $\nabla (R_i^f \vec X_i) = \partial (R_i^f \vec X_i) = 0$ as $R_i^f \vec X_i$ is constant vector field 
inside triangle tangent space (triangle is flat $\mathbb R^2$ space).


\section{Reference Elements}

One of the novelties of FEM is to be able to compute the integrals per individual elements on an easy and generic 
way by using so called Reference Elements.
 
For triangle elements we just need a single reference triangle
 
\begin{center}
	\begin{tikzpicture}
		\draw (0,0) node[anchor=south]{$P_3$}
		-- (0,-3) node[anchor=north]{$P_1$}
		-- (3,-3) node[anchor=north]{$P_2$}
		-- cycle;
	\end{tikzpicture}		
\end{center}
\begin{eqnarray} 
	P_1 = (0,0) \nonumber\\
	P_2 = (1,0) \nonumber\\
	P_3 = (0,1) \nonumber
\end{eqnarray}

The P1 (linear) 2D Lagrange shape functions $\psi(u,v)$ has the form 
\begin{eqnarray} 
	\psi(u,v) = A u + B u + C \nonumber
\end{eqnarray}

For the reference trignale defined above the $\psi(u,v)$ takes the simple form:
\begin{eqnarray} 
	\psi_1(u,v) = 1 - u - v \nonumber\\
	\psi_2(u,v) = u \nonumber\\
	\psi_3(u,v) = v \nonumber
\end{eqnarray}

We can check that $\psi_i \cdot \psi_j = \delta_{ij}$
\begin{eqnarray} 
\psi_1(P_1) = 1 \ \  \psi_1(P_2) = 0  \ \  \psi_1(P_3) = 0 \nonumber\\
\psi_2(P_1) = 0 \ \  \psi_2(P_2) = 1  \ \  \psi_2(P_3) = 0 \nonumber\\
\psi_3(P_1) = 0 \ \  \psi_3(P_2) = 0  \ \  \psi_3(P_3) = 1 \nonumber
\end{eqnarray}

We define a function $F(\vec U)$, a.k.a., \emph{pushforward}, that maps reference 2D triangle to "physical" 2D triangle on the mesh
(the inverse $F^{-1}(\vec X)$ , a.k.a., \emph{pullback}, maps triangle from physical to reference frame).
\begin{eqnarray}
	\left[\begin{array}{c}
		x \\
		y
	\end{array}\right] = 
	\left[\begin{array}{cc}
		x_2 - x_1 &  x_3 - x_1 \\
		y_2 - y_1 &  y_3 - y_1
	\end{array}\right] 
	\left[\begin{array}{c}
		u \\
		v
	\end{array}\right] +
	\left[\begin{array}{c}
		x_1 \\
		y_1
	\end{array}\right] \nonumber
\end{eqnarray}

Which cab be expressed in vector form
\begin{eqnarray} 
	F(\vec U) =  B \vec U + \vec X_1 \nonumber\\
	F^{-1}(\vec X)  = B^{-1} (\vec X - \vec X_1) \nonumber
\end{eqnarray}


We denote $\psi^k_i$ the function $\psi_i$ defined on physical triangle $K$.

\begin{eqnarray} 
	\psi^k_i(\vec X) = \psi_i \circ  F^{-1}(\vec X) \nonumber\\
	\psi^k_i(\vec X) = \psi_i( F^{-1}(\vec X) ) \nonumber
\end{eqnarray}

The gradient of $\psi^k_i(\vec X)$ in physical triangle is:

\begin{eqnarray} 
	\nabla_{\vec X} \psi^k_i(\vec X) = \nabla_{\vec X} \psi_i( F^{-1}(\vec X) ) \nonumber\\
    = \nabla_{\vec U} \psi_i( \vec U ) \cdot \nabla_{\vec X} F^{-1}(\vec X) \nonumber\\
    = B^{-T} \nabla_{\vec U} \psi_i( \vec U) \nonumber
\end{eqnarray}
where we used $\vec U = F^{-1}(\vec X)$ and $\nabla_{\vec X} F^{-1}(\vec X) = B^{-T}$.

The gradient of $\psi_i(\vec U)$ in reference triangle is:

\begin{eqnarray} 
	\nabla \psi_1 = (-1, -1) \nonumber\\
	\nabla \psi_2 = (1, 0) \nonumber\\
	\nabla \psi_3 = (0, 1) \nonumber
\end{eqnarray}

So the gradient of $\psi^k_i(\vec X)$ in physical triangle $k$ is ``constant'':
 
\begin{eqnarray} 
	\nabla \psi^k_1 = B^{-T} (-1, -1) \nonumber\\
	\nabla \psi^k_2 = B^{-T} (1, 0) \nonumber\\
	\nabla \psi^k_3 = B^{-T} (0, 1) \nonumber
\end{eqnarray}

\subsection{For 3D Triangles}

For 3D triangles the map function $F$ changes as follows:
\begin{eqnarray}
	\left[\begin{array}{c}
		x \\
		y \\
		z
	\end{array}\right] = 
	\left[\begin{array}{cc}
		x_2 - x_1 &  x_3 - x_1 \\
		y_2 - y_1 &  y_3 - y_1 \\
		z_2 - z_1 &  z_3 - z_1
	\end{array}\right] 
	\left[\begin{array}{c}
		u \\
		v
	\end{array}\right] +
	\left[\begin{array}{c}
		x_1 \\
		y_1 \\
		z_1
	\end{array}\right] \nonumber
\end{eqnarray}

In matrix algebra it can be expressed as:

\begin{eqnarray}
	F(\vec U) = B \vec U + \vec X_1 \nonumber\\
	F^{-1}(\vec X) = (B^T B)^{-1} B^T (\vec X - \vec X_1)  \nonumber
\end{eqnarray}
or
\begin{eqnarray}
	F^{-1}(\vec X) = B^{-*} (\vec X - \vec X_1) \nonumber
\end{eqnarray}
where $B^{-*} = (B^T B)^{-1} B^T$ is the Moore-Penrose pseudo-inverse of $B$.

The same analysis as before holds so the gradient in ``physical'' trignale $k$ is:

\begin{eqnarray} 
	\nabla \psi^k_1 = B^{-*T} (-1, -1) \nonumber\\
	\nabla \psi^k_2 = B^{-*T} (1, 0) \nonumber\\
	\nabla \psi^k_3 = B^{-*T} (0, 1) \nonumber
\end{eqnarray}
 
\section{Numerical Quadrature}

The entries of the matrix $M_T^k$ are the following double integrals:

\begin{eqnarray} 
	M_T^k(i,j) = \int_{T^k}{ \nabla \psi^k_i(\vec X) \cdot \nabla \psi^k_j(\vec X) \ dA} \nonumber
\end{eqnarray}
inidices $(i,j)$ run for all nodes of triangle $T$.
 
Let define the inner product of basis functions as  

\begin{eqnarray} 
	I^k_{ij}(\vec X) = \nabla \psi^k_i(\vec X) \cdot \nabla \psi^k_j(\vec X) \nonumber
\end{eqnarray}

The 3-point Gauss quadrature rule for the linear triangle is:
\begin{eqnarray} 
	M_T^k(i,j) = \frac{area(T^k)}{3} ( I^k_{ij}(P_1) +  I^k_{ij}(P_2) +  I^k_{ij}(P_3))
\end{eqnarray}
evaluated at the vertices of the triangle $P_1$, $P_2$ and $P_3$. Which is an exact quadrature formula 
for linear function on a triangle. However since $\nabla \psi^k_i(\vec X)$ does not depends on $\vec X$ the quadrature 
expression reduces to:
\begin{eqnarray} 
	M_T^k(i,j) = area(T^k) \nabla \psi^k_i(\vec X) \cdot \nabla \psi^k_j(\vec X) \nonumber
\end{eqnarray}


\section{Implementation}

Implementation can be found in \cite{Basic2DHeatEquationFEM}

\iffalse

\fi

\bibliographystyle{abbrv}
\bibliography{fembasic2d}

% ------------------------------------------------------------------------
\end{document}
% ------------------------------------------------------------------------
